%! Least Squares Approximations — Unit 4, Lesson 4.3 — Guided Notes
\documentclass[10pt]{article}
\usepackage{linalg-article}
\usepackage{linalg-boxes}
\newcommand{\TallMath}[1]{$\displaystyle #1\rule[-1.4em]{0pt}{3.2em}$}
\newcommand{\vv}{\mathbf{v}}
\newcommand{\ww}{\mathbf{w}}
\newcommand{\xx}{\mathbf{x}}
\newcommand{\bb}{\mathbf{b}}
\newcommand{\pp}{\mathbf{p}}
\newcommand{\ee}{\mathbf{e}}
\newcommand{\avec}{\mathbf{a}}
\newcommand{\zero}{\mathbf{0}}
\newcommand{\T}{\mathsf{T}}

\begin{document}

\pageheader{Unit 4, Lesson 4.3}{Guided Notes}
\namedateperiod

\vspace{0.10in}

\begin{objectivebox}
\textbf{By the end of this lesson, I will be able to\dots}
\begin{itemize}\setlength{\itemsep}{2pt}
  \item explain why $A\xx=\bb$ can have \emph{no} solution, and what the ``best'' answer $\hat{\xx}$ means;
  \item fit a straight line $y=C+Dt$ to data by solving the normal equations $A^{\T}A\hat{\xx}=A^{\T}\bb$;
  \item read the fitted values $\pp=A\hat{\xx}$ and the residuals $\ee=\bb-\pp$, and say what least squares makes small.
\end{itemize}
\end{objectivebox}

\vspace{0.04in}

\begin{vocabbox}
\small
Fill in each term as we build it together.
\vspace{2pt}
\termblanklong{Least-squares solution $\hat{\xx}$}
\termblanklong{Best-fit line $y=C+Dt$}
\termblanklong{Residual $\ee=\bb-\pp$}
\end{vocabbox}

\begin{hookbox}
A spring is stretched by hanging weights on it. You measure its length four times, but your ruler is a
little off each time, so the four points do \emph{not} lie on one straight line.
\begin{itemize}\setlength{\itemsep}{1pt}
  \item Can a single straight line pass \emph{exactly} through all four measurements? \writeline
  \item If no line hits every point, what should we mean by the \emph{best} line? \writeline
\end{itemize}
\end{hookbox}

\begin{notesbox}{1. When $A\xx=\bb$ Has No Solution}
Line up four measurements as four equations $C+Dt_i=y_i$ --- two unknowns $C,D$ but \emph{four}
equations. That is the system $A\xx=\bb$ with
\[
  A=\begin{bsmallmatrix} 1 & t_1\\ 1 & t_2\\ 1 & t_3\\ 1 & t_4 \end{bsmallmatrix},\qquad
  \xx=\begin{bsmallmatrix} C\\ D \end{bsmallmatrix},\qquad
  \bb=\begin{bsmallmatrix} y_1\\ y_2\\ y_3\\ y_4 \end{bsmallmatrix}.
\]
Usually $\bb$ is \emph{not} in the column space $C(A)$ (the plane of all $A\xx$), so there is
\blank{2.6cm} exact solution. Instead of giving up, we ask for the point $\pp=A\hat{\xx}$ \emph{in} the
plane that is \textbf{closest} to $\bb$ --- exactly the \textbf{projection} from Lesson~4.2. The best
$\hat{\xx}$ is called the \textbf{least-squares solution}.
\end{notesbox}

\begin{notesbox}{2. ``Closest'' Means the Normal Equations Again}
Closest means the leftover $\ee=\bb-A\hat{\xx}$ is \emph{perpendicular} to every column of $A$, i.e.
$A^{\T}\ee=\zero$. That is the same rearrangement as last lesson:
\[
  A^{\T}(\bb-A\hat{\xx})=\zero \quad\Longrightarrow\quad
  \boxed{\,A^{\T}A\,\hat{\xx}=A^{\T}\bb\,}\quad\text{(the \textbf{normal equations}).}
\]
Solve this ordinary $2\times2$ system for $\hat{\xx}=(C,D)$; then the best-fit line is $y=C+Dt$.
\emph{Least squares $=$ projection onto $C(A)$, written for data.} The name comes from what it makes
small: the total \emph{squared} error $\|\ee\|^2=e_1^2+e_2^2+\cdots$ --- no line makes it smaller.
\end{notesbox}

\begin{notesbox}{3. Worked Example --- Fit a Line to the Spring Data}
A spring's length $y$ (cm) is measured at weights $t=0,1,2,3$ kg: $\bb=(3,3,5,9)$. Fit $y=C+Dt$.
\[
  A=\begin{bsmallmatrix} 1 & 0\\ 1 & 1\\ 1 & 2\\ 1 & 3 \end{bsmallmatrix},\qquad
  A^{\T}A=\begin{bsmallmatrix} 4 & 6\\ 6 & 14 \end{bsmallmatrix},\qquad
  A^{\T}\bb=\begin{bsmallmatrix} 20\\ 40 \end{bsmallmatrix}.
\]
(The top-left $4$ counts the four rows; $6=0{+}1{+}2{+}3$; $14=0{+}1{+}4{+}9$; $20=\sum y$;
$40=\sum t_i y_i$.) Solve the normal equations:
\[
  \begin{bsmallmatrix} 4 & 6\\ 6 & 14 \end{bsmallmatrix}\begin{bsmallmatrix} C\\ D \end{bsmallmatrix}
  =\begin{bsmallmatrix} 20\\ 40 \end{bsmallmatrix}
  \quad\Longrightarrow\quad
  C=\blank{0.8cm},\quad D=\blank{0.8cm}
  \quad\Longrightarrow\quad \hat{\xx}=\begin{bsmallmatrix} 2\\ 2 \end{bsmallmatrix}.
\]
\textbf{Best-fit line:} $y=\blank{1.0cm}+\blank{1.0cm}\,t$, i.e. $y=2+2t$.
\end{notesbox}

\begin{notesbox}{4. Fitted Values, Residuals, and the Picture}
The line's height at each weight is $\pp=A\hat{\xx}=(2,4,6,8)$ --- the \textbf{fitted values}. The
\textbf{residuals} are the vertical misses
\[
  \ee=\bb-\pp=\begin{bsmallmatrix} 3\\3\\5\\9 \end{bsmallmatrix}-\begin{bsmallmatrix} 2\\4\\6\\8 \end{bsmallmatrix}
  =\begin{bsmallmatrix} 1\\-1\\-1\\1 \end{bsmallmatrix}.
\]
They are perpendicular to both columns: $(1,1,1,1)\cdot\ee=\blank{0.8cm}$ and
$(0,1,2,3)\cdot\ee=\blank{0.8cm}$. So $\ee$ lands in the left nullspace $N(A^{\T})$, the orthogonal
complement of $C(A)$ (Lesson~4.1) --- least squares splits the data into a part \emph{on} the line
($\pp$) plus a perpendicular \emph{miss} ($\ee$).

\begin{center}
% Display coords: x = 1.2*t, y = 0.4*(value). Data (t,y): (0,3)(1,3)(2,5)(3,9); line y=2+2t.
\begin{tikzpicture}[scale=1.0, >=stealth]
  \draw[->, charcoal] (-0.3,0)--(4.7,0) node[right, font=\footnotesize]{$t$ (kg)};
  \draw[->, charcoal] (0,-0.3)--(0,4.2) node[above, font=\footnotesize]{$y$ (cm)};
  % best-fit line y=2+2t: t=0 -> (0,0.8), t=3.5 -> (4.2,3.6)
  \draw[burgundy, thick] (0,0.8)--(4.2,3.6) node[right, font=\footnotesize]{$y=2+2t$};
  % residual segments (data point -- line), then data points on top
  \draw[charcoal, dashed] (0,1.2)--(0,0.8);
  \draw[charcoal, dashed] (1.2,1.2)--(1.2,1.6);
  \draw[charcoal, dashed] (2.4,2.0)--(2.4,2.4);
  \draw[charcoal, dashed] (3.6,3.6)--(3.6,3.2);
  \foreach \x/\y in {0/1.2, 1.2/1.2, 2.4/2.0, 3.6/3.6}{
    \fill[royalblue] (\x,\y) circle (2.6pt);
  }
  \node[royalblue, font=\scriptsize] at (3.05,3.55) {data};
\end{tikzpicture}
\end{center}
\end{notesbox}

\begin{practicebox}
Show your work.
\begin{enumerate}[leftmargin=1.6em, itemsep=6pt]
  \item The three points $(0,3),(1,4),(2,5)$ happen to lie \emph{exactly} on one line. What line is it,
        and what are the residuals? (This is the special case where $A\xx=\bb$ \emph{does} have a
        solution.) \writeline
  \item In a least-squares fit, what does a residual $e_i$ measure, and why do we make the sum of their
        \emph{squares} small rather than just their sum? \writeline
\end{enumerate}
\end{practicebox}

\end{document}
